\section{Principal divisors and pushforward}
\label{section-two-fun}

\noindent
The first lemma implies that the pushforward of a principal
divisor along a generically finite morphism is a principal divisor.

\begin{lemma}
\label{lemma-proper-pushforward-alteration}
Let $(S, \delta)$ be as in Situation \ref{situation-setup}.
Let $X$, $Y$ be locally of finite type over $S$. Assume $X$, $Y$
are integral and $n = \dim_\delta(X) = \dim_\delta(Y)$.
Let $p : X \to Y$ be a dominant proper morphism.
Let $f \in R(X)^*$. Set
$$
g = \text{Nm}_{R(X)/R(Y)}(f).
$$
Then we have
$p_*\text{div}(f) = \text{div}(g)$.
\end{lemma}

\begin{proof}
Let $Z \subset Y$ be an integral closed subscheme of $\delta$-dimension
$n - 1$. We want to show that the coefficient of $[Z]$ in
$p_*\text{div}(f)$ and $\text{div}(g)$ are equal. We may apply
Lemma \ref{lemma-finite-in-codimension-one}
to the morphism $p : X \to Y$ and the generic point $\xi \in Z$.
Hence we may replace $Y$ by an
affine open neighbourhood of $\xi$ and assume that $p : X \to Y$ is finite.
Write $Y = \Spec(R)$ and $X = \Spec(A)$ with $p$ induced
by a finite homomorphism $R \to A$ of Noetherian domains which induces
an finite field extension $L/K$ of fraction fields.
Now we have $f \in L$, $g = \text{Nm}(f) \in K$,
and a prime $\mathfrak p \subset R$ with $\dim(R_{\mathfrak p}) = 1$.
The coefficient of $[Z]$ in $\text{div}_Y(g)$ is
$\text{ord}_{R_\mathfrak p}(g)$.
The coefficient of $[Z]$ in $p_*\text{div}_X(f)$ is
$$
\sum\nolimits_{\mathfrak q\text{ lying over }\mathfrak p}
[\kappa(\mathfrak q) : \kappa(\mathfrak p)]
\text{ord}_{A_{\mathfrak q}}(f)
$$
The desired equality therefore follows from
Algebra, Lemma \ref{algebra-lemma-finite-extension-dim-1}.
\end{proof}

\noindent
An important role in the discussion of principal divisors
is played by the ``universal'' principal divisor $[0] - [\infty]$
on $\mathbf{P}^1_S$. To make this more precise, let us denote
\begin{equation}
\label{equation-zero-infty}
D_0, D_\infty \subset
\mathbf{P}^1_S = \underline{\text{Proj}}_S(\mathcal{O}_S[T_0, T_1])
\end{equation}
the closed subscheme cut out by the section $T_1$, resp.\ $T_0$
of $\mathcal{O}(1)$. These are effective Cartier divisors, see
Divisors, Definition \ref{divisors-definition-effective-Cartier-divisor}
and Lemma \ref{divisors-lemma-characterize-OD}.
The following lemma says that loosely speaking we have
``$\text{div}(T_1/T_0) = [D_0] - [D_1]$'' and that this is the
universal principal divisor.

\begin{lemma}
\label{lemma-rational-function}
Let $(S, \delta)$ be as in Situation \ref{situation-setup}.
Let $X$ be locally of finite type over $S$. Assume $X$ is
integral and $n = \dim_\delta(X)$. Let $f \in R(X)^*$.
Let $U \subset X$ be a nonempty open such that $f$
corresponds to a section $f \in \Gamma(U, \mathcal{O}_X^*)$.
Let $Y \subset X \times_S \mathbf{P}^1_S$ be the
closure of the graph of $f : U \to \mathbf{P}^1_S$.
Then
\begin{enumerate}
\item the projection morphism $p : Y \to X$ is proper,
\item $p|_{p^{-1}(U)} : p^{-1}(U) \to U$ is an isomorphism,
\item the pullbacks $Y_0 = q^{-1}D_0$ and $Y_\infty = q^{-1}D_\infty$
via the morphism $q : Y \to \mathbf{P}^1_S$ are defined
(Divisors, Definition
\ref{divisors-definition-pullback-effective-Cartier-divisor}),
\item we have
$$
\text{div}_Y(f) = [Y_0]_{n - 1} - [Y_\infty]_{n - 1}
$$
\item we have
$$
\text{div}_X(f) = p_*\text{div}_Y(f)
$$
\item if we view $Y_0$ and $Y_\infty$ as closed subschemes of $X$
via the morphism $p$ then we have
$$
\text{div}_X(f) = [Y_0]_{n - 1} - [Y_\infty]_{n - 1}
$$
\end{enumerate}
\end{lemma}

\begin{proof}
Since $X$ is integral, we see that $U$ is integral.
Hence $Y$ is integral, and $(1, f)(U) \subset Y$ is an open dense subscheme.
Also, note that the closed subscheme $Y \subset X \times_S \mathbf{P}^1_S$
does not depend on the choice of the open $U$, since after all it is
the closure of the one point set $\{\eta'\} = \{(1, f)(\eta)\}$
where $\eta \in X$ is the generic point. Having said this let us
prove the assertions of the lemma.

\medskip\noindent
For (1) note that $p$ is the composition of the closed immersion
$Y \to X \times_S \mathbf{P}^1_S = \mathbf{P}^1_X$ with the proper
morphism $\mathbf{P}^1_X \to X$. As a composition of proper morphisms
is proper (Morphisms, Lemma \ref{morphisms-lemma-composition-proper})
we conclude.

\medskip\noindent
It is clear that $Y \cap U \times_S \mathbf{P}^1_S = (1, f)(U)$.
Thus (2) follows. It also follows that $\dim_\delta(Y) = n$.

\medskip\noindent
Note that $q(\eta') = f(\eta)$ is not contained in $D_0$ or $D_\infty$
since $f \in R(X)^*$. Hence (3) by
Divisors, Lemma \ref{divisors-lemma-pullback-effective-Cartier-defined}.
We obtain $\dim_\delta(Y_0) = n - 1$
and $\dim_\delta(Y_\infty) = n - 1$ from
Lemma \ref{lemma-divisor-delta-dimension}.

\medskip\noindent
Consider the effective Cartier divisor $Y_0$.
At every point $\xi \in Y_0$ we have $f \in \mathcal{O}_{Y, \xi}$ and
the local equation for $Y_0$ is given by $f$.
In particular, if $\delta(\xi) = n - 1$ so $\xi$ is the generic point
of an integral closed subscheme $Z$ of $\delta$-dimension $n - 1$,
then we see that the coefficient of $[Z]$ in $\text{div}_Y(f)$ is
$$
\text{ord}_Z(f) =
\text{length}_{\mathcal{O}_{Y, \xi}}
(\mathcal{O}_{Y, \xi}/f\mathcal{O}_{Y, \xi}) =
\text{length}_{\mathcal{O}_{Y, \xi}}
(\mathcal{O}_{Y_0, \xi})
$$
which is the coefficient of $[Z]$ in $[Y_0]_{n - 1}$. A similar
argument using the rational function $1/f$ shows that
$-[Y_\infty]$ agrees with the terms with negative coefficients in
the expression for $\text{div}_Y(f)$. Hence (4) follows.

\medskip\noindent
Note that $D_0 \to S$ is an isomorphism. Hence we see that
$X \times_S D_0 \to X$ is an isomorphism as well. Clearly
we have $Y_0 = Y \cap X \times_S D_0$ (scheme theoretic intersection)
inside $X \times_S \mathbf{P}^1_S$. Hence it is really the case that
$Y_0 \to X$ is a closed immersion. It follows that
$$
p_*\mathcal{O}_{Y_0} = \mathcal{O}_{Y'_0}
$$
where $Y'_0 \subset X$ is the image of $Y_0 \to X$.
By Lemma \ref{lemma-cycle-push-sheaf} we
have $p_*[Y_0]_{n - 1} = [Y'_0]_{n - 1}$. The same
is true for $D_\infty$ and $Y_\infty$. Hence (6) is a consequence of (5).
Finally, (5) follows immediately from
Lemma \ref{lemma-proper-pushforward-alteration}.
\end{proof}

\noindent
The following lemma says that the degree of a principal divisor on
a proper curve is zero.

\begin{lemma}
\label{lemma-curve-principal-divisor}
Let $K$ be any field. Let $X$ be a $1$-dimensional integral scheme
endowed with a proper morphism $c : X \to \Spec(K)$.
Let $f \in K(X)^*$ be an invertible rational function.
Then
$$
\sum\nolimits_{x \in X \text{ closed}}
[\kappa(x) : K] \text{ord}_{\mathcal{O}_{X, x}}(f)
=
0
$$
where $\text{ord}$ is as in
Algebra, Definition \ref{algebra-definition-ord}.
In other words, $c_*\text{div}(f) = 0$.
\end{lemma}

\begin{proof}
Consider the diagram
$$
\xymatrix{
Y \ar[r]_p \ar[d]_q & X \ar[d]^c \\
\mathbf{P}^1_K \ar[r]^-{c'} & \Spec(K)
}
$$
that we constructed in Lemma \ref{lemma-rational-function}
starting with $X$ and the rational function $f$ over $S = \Spec(K)$.
We will use all the results of this lemma without further mention.
We have to show that $c_*\text{div}_X(f) = c_*p_*\text{div}_Y(f) = 0$.
This is the same as proving that $c'_*q_*\text{div}_Y(f) = 0$.
If $q(Y)$ is a closed point of $\mathbf{P}^1_K$ then we
see that $\text{div}_X(f) = 0$ and the lemma holds.
Thus we may assume that $q$ is dominant.
Suppose we can show that $q : Y \to \mathbf{P}^1_K$ is finite
locally free of degree $d$ (see
Morphisms, Definition \ref{morphisms-definition-finite-locally-free}).
Since $\text{div}_Y(f) = [q^{-1}D_0]_0 - [q^{-1}D_\infty]_0$
we see (by definition of flat pullback) that
$\text{div}_Y(f) = q^*([D_0]_0 - [D_\infty]_0)$.
Then by Lemma \ref{lemma-finite-flat} we get
$q_*\text{div}_Y(f) = d([D_0]_0 - [D_\infty]_0)$.
Since clearly $c'_*[D_0]_0 = c'_*[D_\infty]_0$ we win.

\medskip\noindent
It remains to show that $q$ is finite locally free.
(It will automatically have some given degree as $\mathbf{P}^1_K$
is connected.)
Since $\dim(\mathbf{P}^1_K) = 1$ we see that $q$ is finite for example
by Lemma \ref{lemma-finite-in-codimension-one}.
All local rings of $\mathbf{P}^1_K$ at
closed points are regular local rings of dimension $1$
(in other words discrete valuation rings), since they are
localizations of $K[T]$ (see
Algebra, Lemma \ref{algebra-lemma-dim-affine-space}).
Hence for $y\in Y$ closed the local ring $\mathcal{O}_{Y, y}$
will be flat over $\mathcal{O}_{\mathbf{P}^1_K, q(y)}$ as soon as
it is torsion free (More on Algebra, Lemma
\ref{more-algebra-lemma-dedekind-torsion-free-flat}).
This is obviously the case as
$\mathcal{O}_{Y, y}$ is a domain and $q$ is dominant.
Thus $q$ is flat. Hence $q$ is finite locally free by
Morphisms, Lemma \ref{morphisms-lemma-finite-flat}.
\end{proof}





